\chapter{网与滤子}
\label{chap:nets-filters-solutions}

\begin{solution}{directed-finite-subsets}
记 \(\mathcal D=\{F\subseteq S:F\text{ 有限}\}\)。包含关系自反、传递，
而 \(F,G\) 的共同上界可取 \(F\cup G\)，故 \(\mathcal D\) 有向。

若 \((F_n)_{n\in\N}\) 是可数子族，则
\(T=\bigcup_{n\in\N}F_n\) 至多可数。若 \(S\) 不可数，可取
\(s\in S\setminus T\)。有限集 \(\{s\}\) 不包含在任何 \(F_n\) 中，
所以 \((F_n)\) 不共尾。因而 \(\mathcal D\) 没有可数共尾子集。
\end{solution}

\begin{solution}{net-eventual-frequent-negation}
“\(P\) 不频繁成立”表示存在 \(i_0\)，使每个 \(i\succcurlyeq i_0\)
都不满足 \(P\)，这正是“\(P\) 最终不成立”。取否命题便得所需等价。

另一方面，“\(P\) 不最终成立”只说明对每个 \(i_0\) 都能在其后找到
一个不满足 \(P\) 的指标，即“非 \(P\) 频繁成立”。它允许 \(P\) 本身也
频繁成立。例如数列指标上令 \(P(n)\) 表示 \(n\) 为偶数，则 \(P\) 与
非 \(P\) 都频繁成立，二者都不最终成立。
\end{solution}

\begin{solution}{subnet-preserves-limit}
设 \(x_i\to x\)，\(y_j=x_{\phi(j)}\) 是子网。给定 \(x\) 的邻域
\(U\)，存在 \(i_0\) 使 \(i\succcurlyeq i_0\) 时 \(x_i\in U\)。
最终共尾性给出 \(j_0\)，使 \(j\succcurlyeq j_0\) 时
\(\phi(j)\succcurlyeq i_0\)。因此这些 \(j\) 都满足
\(y_j=x_{\phi(j)}\in U\)，故 \(y_j\to x\)。
\end{solution}

\begin{solution}{closure-neighborhood-net}
设 \(x\in\overline A\)。令
\[
 I=\{(U,a):U\text{ 是 }x\text{ 的邻域},\ a\in A\cap U\}.
\]
由于每个邻域与 \(A\) 相交，\(I\ne\varnothing\)。规定
\((U,a)\preccurlyeq(V,b)\) 当且仅当 \(V\subseteq U\)。给定两个
指标，\(U\cap V\) 仍是邻域；任取
\(c\in A\cap U\cap V\)，则 \((U\cap V,c)\) 为共同上界。因此
\(I\) 有向。网 \(a_{(U,a)}=a\) 始终在 \(A\) 中。给定邻域 \(W\)，
从任意 \((W,c)\) 开始，所有后继指标的邻域坐标都包含在 \(W\) 中，
相应网值也在 \(W\) 中，所以该网趋于 \(x\)。

若反过来有 \(A\) 中的网趋于 \(x\)，则每个 \(x\) 的邻域都含有该网
的某项，因而与 \(A\) 相交，故 \(x\in\overline A\)。
\end{solution}

\begin{solution}{tail-filter-verification}
记尾部 \(T_{i_0}=\{x_i:i\succcurlyeq i_0\}\)。每个尾部非空。给定
\(i_0,j_0\)，取 \(k\succcurlyeq i_0,j_0\)，则
\(T_k\subseteq T_{i_0}\cap T_{j_0}\)，所以尾部族是滤子基，生成尾滤子。

若 \(x_i\to x\)，则每个邻域 \(U\ni x\) 包含一个尾部，故
\(U\in\mathcal F_x\)，即 \(\mathcal F_x\to x\)。反之，若尾滤子趋于
\(x\)，则每个邻域 \(U\) 属于尾滤子，因而含某个 \(T_{i_0}\)；这正是
\(x_i\to x\)。
\end{solution}

\begin{solution}{filter-base-canonical-net}
令 \(I_{\mathcal B}=\{(B,x):B\in\mathcal B,x\in B\}\)，并规定
\((B,x)\preccurlyeq(C,y)\) 当且仅当 \(C\subseteq B\)。给定
\((B,x),(C,y)\)，由滤子基性质取 \(D\in\mathcal B\) 使
\(D\subseteq B\cap C\)，再取 \(z\in D\)。则 \((D,z)\) 是共同
上界，所以指标集有向。

从指标 \((B,x)\) 开始，所有后继网值都属于 \(B\)；反之，对任意
\(y\in B\)，指标 \((B,y)\) 是其后继，网值就是 \(y\)。故该尾部的
值集恰为 \(B\)。当 \(B\) 遍历滤子基时，全部尾部正好生成
\(\langle\mathcal B\rangle\)。网收敛等价于尾滤子收敛，因而所需
收敛点完全相同，且两者的聚点也相同。
\end{solution}

\begin{solution}{net-cluster-tail-closure}
若 \(x\) 是聚点，给定 \(i_0\) 和邻域 \(U\ni x\)，可取
\(i\succcurlyeq i_0\) 使 \(x_i\in U\)。因此 \(U\) 与尾部
\(T_{i_0}\) 相交，故 \(x\in\overline{T_{i_0}}\)。这对每个
\(i_0\) 成立。

反之，若 \(x\) 属于所有尾部的闭包，则给定邻域 \(U\) 和 \(i_0\)，
由 \(x\in\overline{T_{i_0}}\) 得 \(U\cap T_{i_0}\ne\varnothing\)。
所以存在 \(i\succcurlyeq i_0\) 使 \(x_i\in U\)，即 \(x\) 是聚点。
\end{solution}

\begin{solution}{ultrafilter-cluster-converges}
收敛总蕴含聚集。反过来，设 \(x\) 是超滤子 \(\mathcal U\) 的聚点。
对任一邻域 \(V\ni x\)，若 \(V\notin\mathcal U\)，超滤子的二分性质
给出 \(X\setminus V\in\mathcal U\)。但
\((X\setminus V)\cap V=\varnothing\)，与聚点定义矛盾。故每个
\(V\in\mathcal N(x)\) 都属于 \(\mathcal U\)，即
\(\mathcal U\to x\)。
\end{solution}

\begin{solution}{continuous-filter-image}
由 \(\{f(A):A\in\mathcal F\}\) 生成的滤子包含 \(B\subseteq Y\)，
当且仅当存在 \(A\in\mathcal F\) 使 \(f(A)\subseteq B\)。这等价于
存在 \(A\in\mathcal F\) 使 \(A\subseteq f^{-1}(B)\)。由滤子的向上
封闭性，后者又等价于 \(f^{-1}(B)\in\mathcal F\)。因此两种定义给出
同一滤子。
\end{solution}

\begin{solution}{compact-closed-subspace-filters}
设 \(A\) 是紧致空间 \(X\) 的闭子空间，\(\mathcal F\) 是 \(A\) 上
的滤子。把它提升为 \(X\) 上的滤子
\[
 \widetilde{\mathcal F}
 =\{B\subseteq X:B\cap A\in\mathcal F\}.
\]
紧致性给出聚点 \(x\in X\)。因为 \(A\in\widetilde{\mathcal F}\)，
聚点必在 \(\overline A=A\) 中；把邻域取相对拓扑的形式即可看出 \(x\)
也是原滤子在 \(A\) 中的聚点。因此 \(A\) 紧致。

再设 \(K\) 是 Hausdorff 空间 \(X\) 的紧致子集，并取
\(x\in\overline K\)。集合族
\[
 \mathcal B=\{U\cap K:U\in\mathcal N(x)\}
\]
是 \(K\) 上的滤子基。紧致性使它生成的滤子有聚点 \(k\in K\)。若
\(k\ne x\)，取不交邻域 \(V\ni x\)、\(W\ni k\)。基集
\(V\cap K\) 与 \(W\) 不交，违背 \(k\) 的聚点性质。因此
\(x=k\in K\)，故 \(K\) 闭。
\end{solution}

\begin{solution}{product-coordinate-net-convergence}
若 \(x_i\to x\)，坐标投影 \(\pi_s:X\to X_s\) 连续，故
\(x_i(s)=\pi_s(x_i)\to\pi_s(x)=x(s)\)。

反之，设每个坐标网都收敛。\(x\) 的基本邻域可写成
\[
 U=\bigcap_{k=1}^n\pi_{s_k}^{-1}(U_k),
\]
其中 \(U_k\) 是 \(x(s_k)\) 的邻域。对每个 \(k\)，坐标收敛给出一个
最终条件 \(x_i(s_k)\in U_k\)。有向性把有限多个最终条件合并，得到
\(x_i\in U\) 最终成立。每个邻域都含一个基本邻域，故 \(x_i\to x\)。
\end{solution}

\begin{solution}{first-countable-sequential-closure}
设 \(x\in\overline A\)，并取 \(x\) 的可数邻域基 \((V_n)\)。令
\(U_n=V_0\cap\cdots\cap V_n\)，则 \((U_n)\) 递减且仍为邻域基。
逐项取 \(a_n\in A\cap U_n\)。给定邻域 \(U\ni x\)，存在 \(m\) 使
\(U_m\subseteq U\)，于是 \(n\ge m\) 时 \(a_n\in U_n\subseteq U\)，
故 \(a_n\to x\)。若 \(A\) 对数列极限封闭，则 \(x\in A\)，所以
\(A\) 闭。反向结论对任意空间都成立。

若 \(f:X\to Y\) 保持数列极限而在 \(x\) 不连续，可取
\(f(x)\) 的邻域 \(V\)，使每个 \(x\) 的邻域都含有一点不在
\(f^{-1}(V)\) 中。对上述递减基逐项取
\(x_n\in U_n\setminus f^{-1}(V)\)。则 \(x_n\to x\)，但
\(f(x_n)\) 从不进入 \(V\)，矛盾。因此序列连续蕴含连续；另一方向
由连续性直接得到。
\end{solution}

\begin{solution}{finite-support-product-closure}
乘积拓扑中 \(\boldsymbol1\) 的基本邻域只规定有限坐标
\(F\subseteq S\) 上的值落在各自给定邻域。由于每个因子是离散的
\(\{0,1\}\)，可缩小为要求这些坐标都等于 \(1\)。有限支集点
\(\boldsymbol1_F\) 满足要求，故每个基本邻域与 \(A\) 相交，
\(\boldsymbol1\in\overline A\)。

以有限子集 \(F\) 按包含关系为指标，网 \(a_F=\boldsymbol1_F\) 在每个
固定坐标 \(s\) 上最终等于 \(1\)：从指标 \(\{s\}\) 以后即可。由乘积
空间的坐标判据，\(a_F\to\boldsymbol1\)。

若 \((a_n)\subseteq A\)，令
\(T=\bigcup_n\{s:a_n(s)=1\}\)。每个支集有限，故 \(T\) 至多可数。
不可数性给出 \(s_0\in S\setminus T\)。于是所有 \(a_n(s_0)=0\)，
而 \(\boldsymbol1(s_0)=1\)；坐标投影连续，所以该数列不可能收敛到
\(\boldsymbol1\)。
\end{solution}

\begin{solution}{universal-subnet-ultrafilter}
设原网的尾滤子为 \(\mathcal F_x\)，并取超滤子
\(\mathcal U\supseteq\mathcal F_x\)。令 \(J\) 由三元组
\[
 (i,A,k),\qquad i,k\in I,\quad A\in\mathcal U,\quad
 k\succcurlyeq i,\quad x_k\in A
\]
组成。这样的三元组总存在，因为 \(A\) 与每个尾部的交属于
\(\mathcal U\)，因而非空。规定
\((i,A,k)\preccurlyeq(i',A',k')\) 当且仅当
\(i\preccurlyeq i'\) 且 \(A'\subseteq A\)。给定两个指标，取
\(i'\) 越过两个第一坐标，并在
\(A_1\cap A_2\cap\{x_k:k\succcurlyeq i'\}\) 中取一项 \(x_{k'}\)，
即可得到共同上界。因此 \(J\) 有向。

令 \(\phi(i,A,k)=k\)。第一坐标保证 \(\phi\) 最终共尾，所以
\((x_{\phi(j)})\) 是子网。对任意 \(B\subseteq X\)，超滤子二分性质
给出 \(B\in\mathcal U\) 或 \(X\setminus B\in\mathcal U\)。若
\(B\in\mathcal U\)，从任意第二坐标等于 \(B\) 的指标开始，后继指标
的第二坐标包含于 \(B\)，其网值因而属于 \(B\)。补集情形相同。故该
子网对每个集合都最终落在它或其补集中。
\end{solution}

\begin{solution}{compact-map-closed}
设 \(A\subseteq X\) 闭，则 \(A\) 紧致。若 \(A=\varnothing\)，像闭；
以下设 \(y\in\overline{f(A)}\)，由闭包的网刻画，存在
\((y_i)_{i\in I}\subseteq f(A)\) 趋于 \(y\)。为同时保留所有原像，令
\[
 J=\{(i,a):i\in I,\ a\in A,\ f(a)=y_i\},
\]
并按第一坐标排序。每个 \(y_i\) 都有原像，所以 \(J\) 非空；给定两个
指标，先取第一坐标的共同上界 \(k\)，再取任意 \(a\in A\) 满足
\(f(a)=y_k\)，便得到共同上界 \((k,a)\)。故 \(J\) 有向。

投影 \(p:J\to I\) 最终共尾：给定 \(i_0\)，取一个 \((i_0,a_0)\)，
它之后的所有第一坐标都超过 \(i_0\)。网 \(z_{(i,a)}=a\) 位于紧空间
\(A\) 中，故有一个收敛子网 \(z_{\psi(\ell)}\to a\in A\)。由连续性，
\(f(z_{\psi(\ell)})\to f(a)\)。另一方面，这个像网是原网的子网，
因为 \(p\circ\psi\) 最终共尾，所以它也趋于 \(y\)。Hausdorff空间
中的网极限唯一，给出 \(y=f(a)\in f(A)\)。因此 \(f(A)\) 闭，
\(f\) 为闭映射。

若 \(f\) 还是双射，则逆映射 \(f^{-1}:Y\to X\) 的开集原像为
\(f(U)\)。由于 \(f\) 闭且双射，\(f(U)=Y\setminus f(X\setminus U)\)
为开集，所以 \(f^{-1}\) 连续，\(f\) 是同胚。
\end{solution}
